\documentclass[12pt,a4paper]{article}
\usepackage{fontspec}\setmainfont{Latin Modern Roman}\setsansfont{DejaVu Sans}
\usepackage{amsmath,amssymb,graphicx,geometry,enumitem,longtable,booktabs,fancyhdr,hyperref}
\geometry{left=28mm,right=23mm,top=24mm,bottom=24mm}
\hypersetup{unicode=true,colorlinks=true,linkcolor=black,urlcolor=blue,pdftitle={Giải tích Chương 2 Mục 2.1 Các bài toán tiếp tuyến và vận tốc}}
\setlength{\parindent}{1.2em}\setlength{\parskip}{4pt}\setlength{\headheight}{16pt}
\pagestyle{fancy}\fancyhf{}\fancyhead[L]{\small Chương 2. Giới hạn và đạo hàm}\fancyhead[R]{\small §2.1}\fancyfoot[C]{\thepage}
\newcommand{\fig}[3][.68]{\begin{center}\begin{minipage}{\linewidth}\centering\includegraphics[width=#1\linewidth]{assets/#2.png}\\{\small #3}\end{minipage}\end{center}}
\newcommand{\ex}[1]{\par\medskip\noindent\textbf{Ví dụ #1.}\ }
\newcommand{\sol}{\par\noindent\textit{Lời giải.}\ }
\newcommand{\bt}[1]{\par\smallskip\noindent\textbf{#1.}\ }
\newcommand{\group}[1]{\par\medskip\noindent\textbf{#1}\par}
\begin{document}
\begin{center}
{\Large\bfseries GIẢI TÍCH}\\[5pt]
{\Large\bfseries Chương 2. Giới hạn và đạo hàm}\\[9pt]
{\large\bfseries 2.1. Các bài toán tiếp tuyến và vận tốc}
\end{center}
\noindent\textbf{Nguồn:} James Stewart, Daniel K. Clegg, Saleem Watson, \textit{Calculus: Early Transcendentals}, ấn bản thứ 9 (2020). Phạm vi: phần mở đầu Chương 2 và trọn §2.1, trang in 77–82 (trang PDF 112–117).
\section*{Mở đầu Chương 2}
\fig[.86]{opening}{Ảnh mở đầu chương. Nguồn ảnh trong sách: Icealex / Shutterstock.com.}
Ta biết rằng một vật được thả từ trên cao sẽ rơi ngày càng nhanh. Galileo phát hiện quãng đường vật rơi tỉ lệ với bình phương thời gian đã trôi qua. Giải tích cho phép tính chính xác tốc độ của vật tại bất kỳ thời điểm nào. Bài tập 2.7.11 yêu cầu xác định tốc độ của một người nhảy từ vách đá khi lao xuống biển.

Trong phần \textit{Giới thiệu về giải tích}, ngay trước Chương 1, ta đã thấy ý tưởng giới hạn là nền tảng của các nhánh khác nhau của giải tích. Vì vậy, việc bắt đầu nghiên cứu giải tích bằng cách khảo sát giới hạn và các tính chất của chúng là thích hợp. Loại giới hạn đặc biệt dùng để tìm tiếp tuyến và vận tốc dẫn đến khái niệm trung tâm của phép tính vi phân: \textbf{đạo hàm}.
\section*{2.1. Các bài toán tiếp tuyến và vận tốc}
Trong mục này, ta sẽ thấy giới hạn xuất hiện như thế nào khi tìm tiếp tuyến của một đường cong hoặc vận tốc của một vật.
\subsection*{Bài toán tiếp tuyến}
Từ \textit{tangent} bắt nguồn từ tiếng Latinh \textit{tangens}, nghĩa là ``tiếp xúc''. Ta có thể hình dung tiếp tuyến của một đường cong là đường thẳng tiếp xúc với đường cong và có cùng hướng với đường cong tại tiếp điểm. Làm thế nào diễn đạt chính xác ý tưởng ấy?

Đối với đường tròn, ta có thể theo Euclid, nói rằng tiếp tuyến là đường thẳng $\ell$ cắt đường tròn tại một và chỉ một điểm, như Hình 1(a). Đối với những đường cong phức tạp hơn, định nghĩa này không đủ. Trong Hình 1(b), đường thẳng $\ell$ có vẻ là tiếp tuyến của đường cong $C$ tại $P$, nhưng lại cắt $C$ ở hai điểm.
\fig[.43]{f01}{Hình 1. (a) Tiếp tuyến của đường tròn. (b) Tiếp tuyến của một đường cong có thể cắt đường cong tại điểm khác.}
Cụ thể, xét bài toán tìm tiếp tuyến $\ell$ của parabol $y=x^2$ trong ví dụ sau.
\ex{1} Tìm phương trình tiếp tuyến của parabol $y=x^2$ tại $P(1,1)$.
\sol Ta sẽ tìm được phương trình của tiếp tuyến $\ell$ ngay khi biết hệ số góc $m$. Khó khăn là chỉ biết một điểm $P$ trên $\ell$, trong khi cần hai điểm để tính hệ số góc. Tuy nhiên, có thể tính gần đúng $m$ bằng cách chọn một điểm $Q(x,x^2)$ trên parabol ở gần $P$ (Hình 2), rồi tính hệ số góc $m_{PQ}$ của cát tuyến $PQ$. Từ \textit{secant} bắt nguồn từ tiếng Latinh \textit{secans}, nghĩa là ``cắt''; cát tuyến là đường thẳng cắt một đường cong tại nhiều hơn một điểm.
\fig[.48]{f02}{Hình 2. Tiếp tuyến $\ell$ và cát tuyến $PQ$ của parabol $y=x^2$.}
Chọn $x\ne1$ để $Q\ne P$. Khi đó,
\[m_{PQ}=\frac{x^2-1}{x-1}.\]
Chẳng hạn, với $Q(1.5,2.25)$,
\[m_{PQ}=\frac{2.25-1}{1.5-1}=\frac{1.25}{0.5}=2.5.\]
Hai bảng sau cho $m_{PQ}$ ứng với một số giá trị $x$ gần 1.
\begin{center}\begin{tabular}{rr@{\qquad}rr}\toprule
$x$&$m_{PQ}$&$x$&$m_{PQ}$\\\midrule
2&3&0&1\\1.5&2.5&0.5&1.5\\1.1&2.1&0.9&1.9\\1.01&2.01&0.99&1.99\\1.001&2.001&0.999&1.999\\\bottomrule
\end{tabular}\end{center}
$Q$ càng gần $P$ thì $x$ càng gần 1; theo bảng, $m_{PQ}$ dường như càng gần 2. Điều này gợi ý rằng hệ số góc của tiếp tuyến $\ell$ là $m=2$.

Ta nói hệ số góc của tiếp tuyến là \textbf{giới hạn} của các hệ số góc cát tuyến, và viết bằng ký hiệu
\[\lim_{Q\to P}m_{PQ}=m\qquad\text{và}\qquad\lim_{x\to1}\frac{x^2-1}{x-1}=2.\]
Giả sử hệ số góc tiếp tuyến thực sự bằng 2, dùng dạng phương trình đường thẳng qua một điểm, biết hệ số góc, $y-y_1=m(x-x_1)$ (xem Phụ lục B), ta được tiếp tuyến qua $(1,1)$:
\[y-1=2(x-1)\qquad\text{hay}\qquad y=2x-1.\]\hfill$\square$

Hình 3 minh họa quá trình lấy giới hạn trong Ví dụ 1. Khi $Q$ tiến đến $P$ dọc theo parabol, các cát tuyến tương ứng quay quanh $P$ và tiến đến tiếp tuyến $\ell$.
\fig[.96]{f03}{Hình 3. Hàng trên: $Q$ tiến đến $P$ từ bên phải. Hàng dưới: $Q$ tiến đến $P$ từ bên trái.}
Nhiều hàm xuất hiện trong khoa học không được mô tả bằng phương trình tường minh mà được xác định bởi số liệu thực nghiệm. Ví dụ sau cho thấy cách ước lượng hệ số góc tiếp tuyến của đồ thị một hàm như vậy.
\ex{2} Một laser xung hoạt động bằng cách tích điện trên tụ điện, rồi giải phóng điện tích đột ngột khi phát xung laser. Bảng sau mô tả điện tích $Q$ còn lại trên tụ (đơn vị coulomb) tại thời điểm $t$ (đơn vị giây, kể từ khi phát xung). Dùng số liệu để vẽ đồ thị hàm số và ước lượng hệ số góc tiếp tuyến tại điểm có $t=0.04$. \textit{Lưu ý của sách:} Hệ số góc tiếp tuyến biểu thị dòng điện từ tụ đến laser (đơn vị ampere).
\begin{center}\begin{tabular}{rr}\toprule$t$&$Q$\\\midrule
0&10\\0.02&8.187\\0.04&6.703\\0.06&5.488\\0.08&4.493\\0.1&3.676\\\bottomrule\end{tabular}\end{center}
\sol Trong Hình 4, ta biểu diễn các số liệu đã cho và dùng những điểm ấy vẽ phác một đường cong gần đúng đồ thị của hàm số.
\fig[.70]{f04}{Hình 4. Điện tích $Q$ (coulomb) còn lại theo thời gian $t$ (giây).}
Với $P(0.04,6.703)$ và $R(0,10)$ trên đồ thị, hệ số góc cát tuyến $PR$ là
\[m_{PR}=\frac{10-6.703}{0-0.04}=-82.425.\]
Bảng sau ghi các kết quả tính tương tự đối với những cát tuyến khác.
\begin{center}\begin{tabular}{rr}\toprule$R$&$m_{PR}$\\\midrule
$(0,10)$&$-82.425$\\$(0.02,8.187)$&$-74.200$\\$(0.06,5.488)$&$-60.750$\\$(0.08,4.493)$&$-55.250$\\$(0.1,3.676)$&$-50.450$\\\bottomrule\end{tabular}\end{center}
Từ bảng, ta dự kiến hệ số góc tiếp tuyến tại $t=0.04$ nằm giữa $-74.20$ và $-60.75$. Thật vậy, trung bình cộng hệ số góc của hai cát tuyến gần nhất là
\[\tfrac12(-74.20-60.75)=-67.475.\]
Theo cách này, ta ước lượng hệ số góc tiếp tuyến xấp xỉ $-67.5$.

Một cách khác là vẽ gần đúng tiếp tuyến tại $P$ rồi đo các cạnh của tam giác $ABC$, như Hình 5.
\fig[.73]{f05}{Hình 5. Ước lượng hệ số góc tiếp tuyến bằng tam giác $ABC$; $t$ tính bằng giây, $Q$ tính bằng coulomb.}
Khi đó, hệ số góc tiếp tuyến được ước lượng là
\[-\frac{|AB|}{|BC|}\approx-\frac{8.0-5.4}{0.06-0.02}=-65.0.\]\hfill$\square$
\noindent\textit{Chú thích bên lề của sách:} Ý nghĩa vật lý của kết quả Ví dụ 2 là dòng điện từ tụ đến laser sau $0.04$ giây xấp xỉ $-65$ ampere. \textit{Xem ghi chú về quy ước dấu ở cuối bản dịch.}
\subsection*{Bài toán vận tốc}
Khi quan sát đồng hồ tốc độ của ô tô trong giao thông đô thị, ta thấy tốc độ không giữ nguyên trong thời gian dài; nghĩa là vận tốc ô tô không cố định. Từ việc quan sát đồng hồ, ta cho rằng tại mỗi thời điểm ô tô có một vận tốc xác định; nhưng vận tốc ``tức thời'' được định nghĩa như thế nào?

Xét bài toán vận tốc: tìm vận tốc tức thời của một vật chuyển động thẳng tại một thời điểm cụ thể, biết vị trí của vật ở mọi thời điểm. Trong ví dụ tiếp theo, ta khảo sát vận tốc của một quả bóng rơi. Qua các thí nghiệm cách đây bốn thế kỷ, Galileo phát hiện quãng đường rơi của một vật rơi tự do tỉ lệ với bình phương thời gian rơi. Mô hình này bỏ qua sức cản không khí. Nếu quãng đường rơi sau $t$ giây được ký hiệu $s(t)$ và tính bằng mét, thì gần mặt đất, nhận xét của Galileo được biểu diễn bởi
\[s(t)=4.9t^2.\]
\ex{3} Một quả bóng được thả từ đài quan sát phía trên của tháp CN ở Toronto, cao $450\,\mathrm m$ so với mặt đất. Tìm vận tốc của bóng sau 5 giây.
\fig[.40]{tower}{Tháp CN tại Toronto. Nguồn ảnh trong sách: Steve Allen / Stockbyte / Getty Images.}
\sol Khó khăn khi tìm vận tốc tức thời ở thời điểm 5 giây là ta chỉ xét một thời điểm $t=5$, không có khoảng thời gian nào. Tuy nhiên, có thể tính gần đúng bằng vận tốc trung bình trên khoảng thời gian ngắn $0.1$ giây, từ $t=5$ đến $t=5.1$:
\[\text{vận tốc trung bình}=\frac{\text{độ biến thiên vị trí}}{\text{thời gian đã trôi qua}}.\]
\[\overline v=\frac{s(5.1)-s(5)}{0.1}=\frac{4.9(5.1)^2-4.9(5)^2}{0.1}=49.49\,\mathrm{m/s}.\]
Bảng sau cho kết quả các phép tính tương tự trên những khoảng thời gian ngày càng ngắn.
\begin{center}\begin{tabular}{lr}\toprule Khoảng thời gian&Vận tốc trung bình ($\mathrm{m/s}$)\\\midrule
$5\le t\le5.1$&49.49\\$5\le t\le5.05$&49.245\\$5\le t\le5.01$&49.049\\$5\le t\le5.001$&49.0049\\\bottomrule\end{tabular}\end{center}
Khi rút ngắn khoảng thời gian, vận tốc trung bình dường như càng gần $49\,\mathrm{m/s}$. \textbf{Vận tốc tức thời} tại $t=5$ được định nghĩa là giá trị giới hạn của các vận tốc trung bình trên những khoảng thời gian ngày càng ngắn bắt đầu từ $t=5$. Vì vậy, vận tốc tức thời sau 5 giây dường như bằng $49\,\mathrm{m/s}$.\hfill$\square$

Các phép tính ở bài toán này rất giống những phép tính tìm tiếp tuyến ở đầu mục. Thực vậy, hai bài toán có quan hệ chặt chẽ. Vẽ đồ thị hàm quãng đường rơi của bóng (Hình 6) và xét các điểm $P\bigl(5,4.9(5)^2\bigr)$, $Q\bigl(5+h,4.9(5+h)^2\bigr)$; khi đó hệ số góc cát tuyến $PQ$ là
\[m_{PQ}=\frac{4.9(5+h)^2-4.9(5)^2}{(5+h)-5},\]
đúng bằng vận tốc trung bình trên khoảng thời gian $[5,5+h]$. Vì vậy, vận tốc tại $t=5$ (giới hạn các vận tốc trung bình khi $h$ tiến đến 0) bằng hệ số góc tiếp tuyến tại $P$ (giới hạn các hệ số góc cát tuyến).
\fig[.56]{f06}{Hình 6. Trên: hệ số góc cát tuyến bằng vận tốc trung bình. Dưới: hệ số góc tiếp tuyến bằng vận tốc tức thời.}
Ví dụ 1 và 3 cho thấy muốn giải các bài toán tiếp tuyến và vận tốc, ta phải biết tìm giới hạn. Sau khi nghiên cứu các phương pháp tính giới hạn trong năm mục tiếp theo, ta sẽ trở lại hai bài toán này ở §2.7.
\section*{2.1. Bài tập}
\bt{1} Một bể chứa 1000 gallon nước; nước thoát qua đáy bể trong nửa giờ. Bảng cho thể tích $V$ còn lại (gallon) sau $t$ phút.
\begin{center}\begin{tabular}{lrrrrrr}\toprule
$t$ (phút)&5&10&15&20&25&30\\
$V$ (gallon)&694&444&250&111&28&0\\\bottomrule\end{tabular}\end{center}
(a) Với $P(15,250)$ trên đồ thị $V$, tìm hệ số góc các cát tuyến $PQ$ khi $Q$ ứng với $t=5,10,20,25,30$.

(b) Ước lượng hệ số góc tiếp tuyến tại $P$ bằng trung bình cộng hệ số góc của hai cát tuyến.

(c) Dùng đồ thị $V$ để ước lượng hệ số góc tiếp tuyến tại $P$. Hệ số này biểu thị tốc độ nước chảy ra khỏi bể sau 15 phút.
\bt{2} Một sinh viên mua đồng hồ thông minh ghi số bước đi trong ngày. Bảng cho số bước được ghi nhận sau $t$ phút kể từ 3 giờ chiều trong ngày đầu đeo đồng hồ.
\begin{center}\begin{tabular}{lrrrrr}\toprule
$t$ (phút)&0&10&20&30&40\\
Số bước&3438&4559&5622&6536&7398\\\bottomrule\end{tabular}\end{center}
(a) Tìm hệ số góc các cát tuyến ứng với (i) $[0,40]$, (ii) $[10,20]$, (iii) $[20,30]$. Các hệ số góc này biểu thị điều gì?

(b) Ước lượng nhịp đi bộ, tính bằng bước/phút, lúc 3 giờ 20 phút chiều bằng trung bình cộng hệ số góc của hai cát tuyến.
\bt{3} Điểm $P(2,-1)$ thuộc đường cong $y=1/(1-x)$.

(a) Với $Q\bigl(x,1/(1-x)\bigr)$, tìm hệ số góc cát tuyến $PQ$, làm tròn đến sáu chữ số sau dấu phẩy, khi:
\[\begin{gathered}(i)\ 1.5;\ (ii)\ 1.9;\ (iii)\ 1.99;\ (iv)\ 1.999;\\[3pt](v)\ 2.5;\ (vi)\ 2.1;\ (vii)\ 2.01;\ (viii)\ 2.001.\end{gathered}\]
(b) Dựa vào câu (a), dự đoán hệ số góc tiếp tuyến tại $P(2,-1)$.

(c) Dùng hệ số góc ở câu (b) tìm phương trình tiếp tuyến tại $P(2,-1)$.
\bt{4} Điểm $P(0.5,0)$ thuộc đường cong $y=\cos\pi x$.

(a) Với $Q(x,\cos\pi x)$, tìm hệ số góc cát tuyến $PQ$, làm tròn đến sáu chữ số sau dấu phẩy, khi:
\[\begin{gathered}(i)\ 0;\ (ii)\ 0.4;\ (iii)\ 0.49;\ (iv)\ 0.499;\\[3pt](v)\ 1;\ (vi)\ 0.6;\ (vii)\ 0.51;\ (viii)\ 0.501.\end{gathered}\]
(b) Dựa vào câu (a), dự đoán hệ số góc tiếp tuyến tại $P(0.5,0)$.

(c) Dùng hệ số góc ở câu (b) tìm phương trình tiếp tuyến tại $P(0.5,0)$.

(d) Vẽ phác đường cong, hai cát tuyến và tiếp tuyến.
\bt{5} Mặt cầu cao 275 feet so với sông. Một viên sỏi rơi khỏi mép cầu; độ cao của nó so với mặt nước sau $t$ giây, tính bằng feet, là $y=275-16t^2$.

(a) Tìm vận tốc trung bình trên khoảng thời gian bắt đầu tại $t=4$ và kéo dài: (i) $0.1$ giây; (ii) $0.05$ giây; (iii) $0.01$ giây.

(b) Ước lượng vận tốc tức thời sau 4 giây.
\bt{6} Một hòn đá được ném thẳng đứng lên trên tại sao Hỏa với vận tốc $10\,\mathrm{m/s}$. Độ cao sau $t$ giây, tính bằng mét, là $y=10t-1.86t^2$.

(a) Tìm vận tốc trung bình trên: (i) $[1,2]$; (ii) $[1,1.5]$; (iii) $[1,1.1]$; (iv) $[1,1.01]$; (v) $[1,1.001]$.

(b) Ước lượng vận tốc tức thời tại $t=1$.
\bt{7} Bảng cho vị trí của một người đi mô tô sau khi bắt đầu tăng tốc từ trạng thái đứng yên.
\begin{center}\begin{tabular}{lrrrrrrr}\toprule
$t$ (giây)&0&1&2&3&4&5&6\\
$s$ (feet)&0&4.9&20.6&46.5&79.2&124.8&176.7\\\bottomrule\end{tabular}\end{center}
(a) Tìm vận tốc trung bình trên: (i) $[2,4]$; (ii) $[3,4]$; (iii) $[4,5]$; (iv) $[4,6]$.

(b) Dùng đồ thị $s$ theo $t$ để ước lượng vận tốc tức thời tại $t=3$.
\bt{8} Độ dời, tính bằng centimét, của một chất điểm chuyển động qua lại trên đường thẳng được cho bởi phương trình chuyển động $s=2\sin\pi t+3\cos\pi t$, trong đó $t$ tính bằng giây.

(a) Tìm vận tốc trung bình trên: (i) $[1,2]$; (ii) $[1,1.1]$; (iii) $[1,1.01]$; (iv) $[1,1.001]$.

(b) Ước lượng vận tốc tức thời của chất điểm tại $t=1$.
\bt{9} Điểm $P(1,0)$ thuộc đường cong $y=\sin(10\pi/x)$.

(a) Với $Q\bigl(x,\sin(10\pi/x)\bigr)$, tìm hệ số góc cát tuyến $PQ$, làm tròn đến bốn chữ số sau dấu phẩy, khi $x=2,1.5,1.4,1.3,1.2,1.1,0.5,0.6,0.7,0.8,0.9$. Các hệ số góc có vẻ tiến đến một giới hạn không?

(b) Dùng đồ thị đường cong để giải thích vì sao các hệ số góc ở câu (a) không gần hệ số góc tiếp tuyến tại $P$. \textit{(Dùng công cụ vẽ đồ thị.)}

(c) Chọn những cát tuyến thích hợp để ước lượng hệ số góc tiếp tuyến tại $P$.
\section*{Ghi chú biên tập}
\noindent\textbf{Quy ước dấu trong Ví dụ 2.} $Q$ là điện tích còn lại trên tụ nên $dQ/dt<0$. Các ước lượng $-67.5$ và $-65.0\,\mathrm A$ là tốc độ biến thiên có dấu của $Q$. Nếu chọn chiều dương của dòng điện từ tụ sang laser, thì $i_{\mathrm{ra}}=-dQ/dt$, có độ lớn tương ứng xấp xỉ $67.5$ hoặc $65.0\,\mathrm A$. Chú thích gốc được giữ nguyên và cần đọc cùng quy ước chiều dòng điện.

\noindent\textbf{Chiều dương trong Ví dụ 3.} $s(t)$ đo quãng đường rơi từ điểm thả xuống dưới, nên chiều dương hướng xuống; $49\,\mathrm{m/s}$ là vận tốc hướng xuống. Nếu chọn trục thẳng đứng hướng lên thì vận tốc có dấu âm. Phân biệt \textit{tốc độ} (đại lượng không âm) với \textit{vận tốc} có dấu trong chuyển động thẳng.

\noindent\textbf{Hình và đơn vị.} Giữ các đường cong, điểm, nhãn ký hiệu của hình gốc; chú thích tiếng Việt giải nghĩa những nhãn tiếng Anh. Đơn vị và dấu chấm thập phân được giữ như sách. Bài tập chỉ được dịch đề, không kèm đáp án.
\section*{Thuật ngữ bổ sung}
\begin{center}\begin{tabular}{ll}\toprule Tiếng Anh&Tiếng Việt thống nhất\\\midrule
limit&giới hạn\\
derivative&đạo hàm\\
differential calculus&phép tính vi phân\\
tangent line&tiếp tuyến\\
point of contact&tiếp điểm\\
secant line&cát tuyến\\
limiting process&quá trình lấy giới hạn\\
average velocity&vận tốc trung bình\\
instantaneous velocity&vận tốc tức thời\\
speed&tốc độ\\
displacement&độ dời\\
free fall&rơi tự do\\
air resistance&sức cản không khí\\
capacitor&tụ điện\\
electric charge&điện tích\\
electric current&dòng điện\\
pulse laser&laser xung\\
\bottomrule\end{tabular}\end{center}
\noindent\textbf{Phạm vi hoàn thành:} Phần mở đầu Chương 2 và trọn §2.1, gồm 3 ví dụ, Hình 1–6, 2 ảnh và 9 bài tập. Phần tiếp theo: \textit{§2.2 — Giới hạn của hàm số}, bắt đầu ở trang in 83.
\end{document}
